%%%%%%%%%%%%%%%%%%%%% chapter.tex %%%%%%%%%%%%%%%%%%%%%%%%%%%%%%%%%
%
% sample chapter
%
% Use this file as a template for your own input.
%
%%%%%%%%%%%%%%%%%%%%%%%% Springer Nature %%%%%%%%%%%%%%%%%%%%%%%%%%
\motto{Imagination is more important than knowledge. Knowledge is limited.
Imagination encircles the world. ---Albert Einstein}
\chapter{The equations}
\label{equations}

\abstract*{
  This chapter will be all about the two equations of the general 
  theory of relativity: the field equations and the geodesic equation.
  With these equations all of the gravitational effects can be explained
  and calculated. It is the field equations that are the primary 
  accomplishment of Einstein's initial work on this theory but the 
  solutions and the solving of the geodesic equation will play a much 
  bigger role in the rest of the book.
}

\abstract{
  This chapter will be all about the two equations of the general 
  theory of relativity: the field equations and the geodesic equation.
  With these equations all of the gravitational effects can be explained
  and calculated. It is the field equations that are the primary 
  accomplishment of Einstein's initial work on this theory but the 
  solutions and the solving of the geodesic equation will play a much 
  bigger role in the rest of the book.
}

\section{The field equations}
\label{sec:field-equations}

\subsection{Einstein Tensor}
\label{subsec:einstein-tensorImagination is more important than knowledge. Knowledge is limited.
Imagination encircles the world. ---Albert Einstein}

\subsection{Derivation}
\label{subsec:field-equations-derivation}

\section{The geodesic equation}
\label{sec:geodesic-equation}

%%%%%%%%%%%%%%%%%%%%%%%%%%%%%%%%%%%%%%%%%%%%%%%%%%%%%%%%%%%%%%%%%%%%%%%%%%%
% References
%%%%%%%%%%%%%%%%%%%%%%%%%%%%%%%%%%%%%%%%%%%%%%%%%%%%%%%%%%%%%%%%%%%%%%%%%%%

\bibliographystyle{spmpsci}
\bibliography{ref-equations}
